\section{The interop budget and the design law}
\label{sec:budget}

We now fold the theorem (Section~\ref{sec:theorem}) and the real anchor
(Section~\ref{sec:interephemeris}) into a multi-provider interoperability budget and
its consistency-tolerance map. Every quantity in this section is \Modelled{} and
tagged \modtag{}: the ``providers'' are the three ephemerides, a representative
analog for a real multi-provider constellation---\emph{no two lunar-PNT providers
fly}---so the budget totals and tolerances describe the structure and ordering of the
trade, not a certified system budget.

\subsection{The budget under three frame conventions}
\label{sec:budgetconventions}

Combining the three provider pairs of Table~\ref{tab:anchor} gives a
representative cross-provider position-inconsistency budget, evaluated under three
interoperability conventions of increasing strength: (i) \code{PerProvider}, each
user keeps each provider in its own realized frame and dynamics; (ii)
\code{CommonFrameTie}, users adopt a common frame realization, removing the reducible
frame-tie $\reduc$ but not the irreducible dynamics $\irr$---this is what today's
format-and-frame standards achieve; and (iii) \code{CommonEphemeris}, users adopt a
common designated dynamical reference, removing $\irr$ as well. Table~\ref{tab:budget}
reports the result. Following the paper's convention, the rotations are the primary
quantities and metres are envelopes: the reducible and irreducible columns are the
RMS-combined frame-tie ($\SI{0.8182}{m}$, from
$|\rot_{\mathrm{tie}}|_{\mathrm{rms}}=\SI{2.131}{nrad}$) and Moon-excess
($\SI{1.7931}{m}$, from $|\rot_{\mathrm{exc}}|_{\mathrm{rms}}=\SI{4.669}{nrad}$)
\emph{worst-case envelopes} $|\rot|\,\Rmoon$ at the Earth--Moon lever arm. The
\code{PerProvider} total is the \emph{measured raw RMS} disagreement
($\SI{1.855}{m}$), not the sum of the two envelopes, because $\rot_{\mathrm{tie}}$ and
$\rot_{\mathrm{exc}}$ are non-parallel and the envelopes do not add in quadrature; each
convention then zeroes the components it can remove. One further caveat is \Modelled{}:
the three pairwise differences are linearly dependent
($(A{-}B)+(B{-}C)=(A{-}C)$), so RMS-combining them overcounts the independent
information---a deliberate, stated modelling choice for a representative envelope.

\begin{table}[t]
\centering
\caption{Multi-provider interop budget (\Modelled{}, \modtag{}; the ``providers''
are the three ephemerides, a representative analog). The reducible and irreducible
columns are RMS-combined worst-case lever-arm envelopes $|\rot|\,\Rmoon$ at
$\Rmoon=\SI{3.840e8}{m}$ (upper bounds; non-additive, since $\rot_{\mathrm{tie}}$ and
$\rot_{\mathrm{exc}}$ are non-parallel). The \code{PerProvider} total is the
\emph{measured raw RMS} disagreement ($\SI{1.855}{m}$), not a quadrature sum;
\code{CommonFrameTie} removes the reducible frame-tie and leaves an
irreducible-dominated $\approx\!\SI{1.8}{m}$ geocentric-Moon-state floor; only
\code{CommonEphemeris} zeroes it. The irreducible fraction
$\|\rot_{\mathrm{exc}}\|/(\|\rot_{\mathrm{tie}}\|+\|\rot_{\mathrm{exc}}\|)=0.687$
(a linear ratio of the RMS rotations) is convention-independent. The three pairwise
differences are linearly dependent ($(A{-}B)+(B{-}C)=(A{-}C)$), so the RMS combination
overcounts---a \Modelled{} choice.}
\label{tab:budget}
\small
\begin{tabular}{@{}lcccc@{}}
\toprule
\textbf{Convention} & \textbf{reducible env.\ (m)} & \textbf{irreducible env.\ (m)} &
\textbf{total (m)} & \textbf{irreducible fraction} \\
\midrule
\code{PerProvider}    & $0.8182$ & $1.7931$ & $1.8550$ & $0.6867$ \\
\code{CommonFrameTie} & $0.8182$ & $1.7931$ & $1.8062$ & $0.6867$ \\
\code{CommonEphemeris}& $0.8182$ & $1.7931$ & $0.0000$ & $0.6867$ \\
\bottomrule
\end{tabular}
\end{table}

The design law follows directly. The irreducible fraction is $0.687$: the irreducible
Moon-excess rotation is about twice the reducible frame-tie, so more than two-thirds of
the cross-provider orientation budget lives in the part that no coordinate convention
removes. Adopting a common frame tie---agreeing on frame tags, format, and an ICRF
realization---lowers the geocentric-Moon-state budget only from the measured
$\SI{1.855}{m}$ to $\SI{1.8062}{m}$, leaving a $\approx\!\SI{1.8}{m}$
irreducible-dominated floor (a Moon-state inconsistency at the Earth--Moon lever arm;
the same excess rotation is only $\approx\!\SI{8}{mm}$ for a lunar-surface user); only
a common designated dynamical reference drives the total to zero. The budget is
monotone in convention strength,
\code{CommonEphemeris}\,$\le$\,\code{CommonFrameTie}\,$<$\,\code{PerProvider}, as
Corollary~\ref{cor:designlaw} requires. This is the concrete, real-data-backed input
to an interoperability specification: \emph{a common format and frame realization is
necessary but not sufficient; the cross-provider floor is set by the lunar-orbit
dynamics, and closing it requires mandating a common (or consistency-bounded)
ephemeris.}

\subsection{The consistency tolerance map}
\label{sec:tolmap}

Inverting Theorem~\ref{thm:tolerance} for a lunar-surface user at
$\Rmoon=R_{\mathrm{Moon}}=\SI{1737400}{m}$, with the per-provider realization taken
as exact ($B_{\mathrm{eff}}=B$), gives the consistency tolerances of
Table~\ref{tab:tolerance}: the maximum inter-provider disagreement, per Helmert
parameter, that keeps a mixed-provider surface user within a position budget $B$. Note
the lever arm here is the lunar \emph{surface} radius $R_{\mathrm{Moon}}$, two orders
of magnitude smaller than the Earth--Moon lever arm at which the
$\approx\!\SI{1.8}{m}$ geocentric-Moon-state floor of Table~\ref{tab:budget} is
evaluated; the two tables therefore speak to different users (a surface asset versus an
ephemeris/cislunar hand-off) and must not be compared metre-for-metre.

\begin{table}[t]
\centering
\caption{Consistency tolerance $\tau(B)$ (\Modelled{}, \modtag{}) for a lunar-surface
user at $\Rmoon=\SI{1737400}{m}$, from the triangle bound of
Theorem~\ref{thm:tolerance} with $B_{\mathrm{eff}}=B$. The origin tolerance is $B$;
the scale and rotation tolerances coincide at $B/\Rmoon$ and are the binding (tightest)
constraints---an orientation error is the cheapest way to exhaust the budget at the
lunar lever arm.}
\label{tab:tolerance}
\small
\begin{tabular}{@{}lcccc@{}}
\toprule
$\bm{B}$ \textbf{(m)} & \textbf{max origin (m)} & \textbf{max scale} &
\textbf{max rotation (rad)} & \textbf{max rotation (nrad)} \\
\midrule
$5.0$  & $5.0000$  & $2.8779\times10^{-6}$ & $2.877863\times10^{-6}$ & $2877.86$ \\
$10.0$ & $10.0000$ & $5.7557\times10^{-6}$ & $5.755727\times10^{-6}$ & $5755.73$ \\
$15.0$ & $15.0000$ & $8.6336\times10^{-6}$ & $8.633590\times10^{-6}$ & $8633.59$ \\
\bottomrule
\end{tabular}
\end{table}

For every budget the binding constraint is the rotation (equivalently the scale): at
the lunar radius a rotation of $\SI{2877.86}{\nano\radian}$ already displaces a
surface asset by the full $\SI{5}{m}$, whereas the origin can drift the full
$\SI{5}{m}$ before it binds. This is consistent with, and explains the practical
importance of, the real anchor being orientation-dominated
(Section~\ref{sec:interephemeris}): the disagreement that actually exists between
lunar providers is exactly of the kind---orientation---to which a surface user's
budget is most sensitive. A specification writer can read Table~\ref{tab:tolerance} as
the required cross-provider rotation agreement: to guarantee a $\SI{5}{m}$
mixed-provider surface budget, compliant providers must agree in orientation to
better than $\approx\!\SI{2.9}{\micro\radian}$. Whether they do so today is the
question the anchor answers negatively for the irreducible Moon-excess: the measured
excess rotations ($1.5$--$6.3$~nrad, Table~\ref{tab:anchor}) are nearly three orders of
magnitude \emph{below} the $\SI{2877.86}{\nano\radian}$ tolerance (a factor of
$460$--$1900$), so at a $\SI{5}{m}$ budget the current ephemeris-level orientation
disagreement is not itself budget-threatening; the map's value is that it sets the
tolerance a real provider constellation must not exceed, and identifies orientation as
the parameter to watch.

\subsection{Application to LNIS AD-5 and multi-provider procurement}
\label{sec:application}

Two immediate uses follow, both stated as \Modelled{} guidance rather than a
certified budget. First, Table~\ref{tab:tolerance} supplies the functional form and
representative magnitude of the cross-provider consistency tolerance that AD-5---the
still-unreleased Lunar Reference System and LunaNet Reference Time System Standard
({LNIS-TBD-AD0005})---does not yet specify: a user PVT budget $B$ maps to a
per-parameter cross-provider tolerance, binding on orientation,
$\tau_\theta=B/R_{\mathrm{Moon}}$. The alignment is pointed: the scope AD-5 itself
declares---a lunar gravity field to finite degree and order together with a parametric
precession/nutation model---is exactly the lunar-orbit dynamics that our anchor
identifies as the irreducible floor, so a standard at the AD-5 level is already scoped
over the very term that dominates the budget. Second, Table~\ref{tab:budget} gives a
vendor-neutral design law for a multi-provider interoperability framework such as
LunaNet, or an ESA procurement such as Moonlight/LCNS: mandating a common message
format and frame realization is necessary but removes only the sub-metre reducible
part; a residual floor set by the providers' independent lunar-orbit dynamics remains,
and a programme that seeks metre-consistent multi-provider PNT must additionally
require a common designated ephemeris, or bound the permissible ephemeris divergence,
as an explicit interoperability criterion. We reiterate the honesty boundary: the
ephemerides stand in for providers, the budget totals and tolerances are
representative, and the reducible/irreducible attribution carries the caveat of
Remark~\ref{rem:attribution}. The convention-free content is robust: an irreducible,
orientation-type, Moon-specific disagreement exists between the real lunar dynamical
references, and no coordinate convention removes it.
